% Abstract for the Introspective Verification paper.
% Style: plain, direct; no em-dashes; minimal colons/quotes; no anthropomorphic verbs.
% Numbers verified against GenPRM (arXiv:2504.00891) Table 1 and ProcessBench.
% Citation keys used elsewhere: processbench, genprm, mathshepherd, azaria2023internal.

Process reward models score the correctness of individual reasoning steps, and the most accurate models today produce this score by generating a natural language critique and, in some systems, by executing code before a verdict. This procedure is accurate but expensive, because every step requires autoregressive decoding and an external interpreter. We ask whether the same judgment is available without generation. We find that step correctness is already encoded in the hidden states of the verifier backbone, and that an intermediate layer near seventy percent of the depth separates correct from incorrect steps more strongly than the final layer when probed. We call the resulting operation introspection: the verdict for a step is decoded from the backbone's own intermediate representation instead of being generated. This probed signal is stronger at the intermediate layer because the final state is optimized to emit the verdict token while an earlier state is formed before that token is produced. Once a readout is trained the final layer becomes the stronger single view, and the residual head combines the two. We build an introspective verdict head (LVH), a small residual classifier that combines one intermediate layer with the final layer and returns a per step reward in a single forward pass, with no generation and no code execution. The backbone is adapted with LoRA on step labels, and both training and inference run on a single V100 at the 1.5B and the 7B scale. On ProcessBench with a 1.5B backbone, the head reaches 58.1 mean F1. This exceeds most process reward models on the same benchmark, including trained process reward models at 7B and 14B scale, and it approaches the 61.9 of GPT-4o, a much larger general model used as a critic without a process reward model, while processing about one quarter of the tokens of a single generative pass. The head is also complementary to generative verifiers. Combined with a 7B generative model in logit space, on 120 case slices and with an unsupervised per domain calibration, it raises mean F1 from 76.8 to 77.4 without lowering any subset.
